% =====================================================================
%  Lecture 05 -- 4G LTE Architecture                      (English edition)
% =====================================================================
\documentclass[11pt,a4paper]{article}
\providecommand\ACNroot{../..}
\input{\ACNroot/templates/lecture-en}

\ACNset{number}{05}
\ACNset{title}{4G LTE Architecture}
\ACNset{shorttitle}{4G LTE Architecture}
\ACNset{titleother}{معماری LTE نسل چهارم}
\ACNset{part}{Part II -- 4G access and core (Lectures 5--13)}
\ACNset{time}{6--8 hours (3 h reading, 3--5 h exercises)}
\ACNset{level}{Intermediate}
\ACNset{prereq}{Lectures 1--4; OFDM basics; IP networking}
\ACNset{gen}{4}

\begin{document}
\ACNcover

% ---------------------------------------------------------------------
\section{Learning objectives}
After studying this lecture you should be able to:
\begin{objectives}
  \item \textbf{Explain} the limitations of 3G that motivated LTE and list the main LTE
        requirements.
  \item \textbf{State} the LTE design principles -- flat, all-IP, shared channels, no soft
        handover -- and the engineering reason for each.
  \item \textbf{Draw} the Evolved Packet System (E-UTRAN + EPC) and \textbf{name} every
        reference point (S1-MME, S1-U, X2, S5/S8, S6a, S11, SGi, Gx, Rx).
  \item \textbf{Derive} the LTE downlink peak rate from the resource grid and
        \textbf{build} a user-plane latency budget.
  \item \textbf{Describe} how voice is carried in an all-IP system (IMS, VoLTE, CS fallback,
        SRVCC).
  \item \textbf{Explain} control/user plane separation (CUPS) and its role as a bridge to 5G.
\end{objectives}

% ---------------------------------------------------------------------
\section{Prerequisite knowledge}
You should know the UMTS architecture and its CS/PS split (Lecture~3), the GTP tunnel of
GPRS (Lecture~1) and the basic idea of OFDM: a wideband channel divided into many narrow,
orthogonal subcarriers, each protected from inter-symbol interference by a cyclic prefix.

\paragraph{Key terms} The vocabulary of this lecture with its Persian equivalents:
\begin{keyterms}
LTE & \fa{تکامل بلندمدت} & 3GPP radio access of 4G (E-UTRA), from Release 8. \\
EPS & \fa{سامانهٔ بستهٔ تکامل‌یافته} & Complete 4G system: E-UTRAN + EPC. \\
eNB (eNodeB) & \fa{ایستگاه پایهٔ LTE} & Base station incl.\ radio resource management; no RNC. \\
EPC & \fa{هستهٔ بستهٔ تکامل‌یافته} & All-IP core: MME, S-GW, P-GW, HSS, PCRF. \\
MME & \fa{موجودیت مدیریت تحرک} & Control-plane node: NAS, mobility, security, bearer control. \\
S-GW / P-GW & \fa{دروازهٔ خدمت‌دهنده / دروازهٔ PDN} & User-plane anchors: mobility anchor / IP anchor. \\
EPS bearer & \fa{حامل EPS} & End-to-end QoS pipe between UE and P-GW. \\
OFDMA / SC-FDMA & \fa{دسترسی چندگانهٔ متعامد / تک‌حامل} & LTE downlink / uplink multiple access. \\
VoLTE & \fa{صدا روی LTE} & IMS-based voice over the EPS bearer of QCI 1. \\
CUPS & \fa{جدایی صفحهٔ کنترل و کاربر} & Rel-14 split of S/P-GW into C and U parts (PFCP). \\
\end{keyterms}

% ---------------------------------------------------------------------
\section{Introduction and motivation}
By 2004 HSPA was delivering mobile broadband, but its foundations were showing their age.
WCDMA's 5\,MHz carrier could not grow to the 10--20\,MHz needed for higher rates without
complex equalisers; the hierarchy UE--NodeB--RNC--SGSN--GGSN added latency and cost; two
parallel cores (CS and PS) doubled operations. Meanwhile competing technologies (WiMAX)
were promising OFDM-based, all-IP broadband.

3GPP's Long Term Evolution study set ambitious targets \cite{tr25913}: peak rates of
100\,Mbit/s downlink and 50\,Mbit/s uplink in 20\,MHz; user-plane latency below 5\,ms;
idle-to-active transition below 100\,ms; 2--4 times the spectral efficiency of Release-6
HSPA; scalable bandwidth from 1.4 to 20\,MHz; and support for speeds up to 350\,km/h. The
\emph{System Architecture Evolution} (SAE) defined a matching core, the EPC. Together they
form the Evolved Packet System (EPS), first specified in Release~8 (2008)
\cite{ts36300,ts23401}.

% ---------------------------------------------------------------------
\section{Fundamental concepts}

\subsection{Design principles and their reasons}
\begin{itemize}
  \item \textbf{Flat RAN.} The RNC is removed; its functions move to the eNB (radio resource
        control, scheduling, ciphering, handover decision) or to the core. Fewer nodes means
        lower latency and cost. The price is that no node can combine signals from several
        cells -- hence no soft handover (Lecture~4).
  \item \textbf{All-packet core.} There is no circuit-switched domain. Voice must be carried as
        IP packets (VoLTE over IMS), or temporarily handed to 2G/3G (CS fallback).
  \item \textbf{Shared channels only.} All user data uses shared channels scheduled every 1\,ms
        by the eNB; there are no dedicated traffic channels.
  \item \textbf{OFDMA downlink, SC-FDMA uplink.} OFDM scales gracefully with bandwidth and makes
        frequency-selective scheduling and MIMO simple. In the uplink, single-carrier FDMA keeps
        a low peak-to-average power ratio for the handset's amplifier (Lecture~7).
  \item \textbf{Separation of control and user planes} in the core: the MME never touches user
        packets; the gateways never terminate NAS signalling.
  \item \textbf{Self-organising networks (SON).} Automatic neighbour relations, mobility robustness
        and load balancing reduce manual planning in a network of many small, flat nodes.
\end{itemize}

\subsection{The EPS bearer}
\begin{definitionbox}{EPS bearer}
An EPS bearer is a logical pipe between the UE and the P-GW with a single QoS treatment. It is
the concatenation of a \emph{radio bearer} (UE--eNB), an \emph{S1 bearer} (eNB--S-GW, a GTP-U
tunnel) and an \emph{S5/S8 bearer} (S-GW--P-GW, another GTP-U tunnel). Every UE gets a
\emph{default bearer} at attach; \emph{dedicated bearers} are added for traffic needing
different QoS, such as VoLTE (QCI~1).
\end{definitionbox}
Unlike GPRS, where connectivity required a separate PDP context activation after attach, LTE
establishes the default bearer \emph{during} attach: an LTE device is always on with an IP
address. Bearers and QoS classes are developed in Lectures~9 and~22.

% ---------------------------------------------------------------------
\section{Detailed technical explanation}

\subsection{Network elements}
\paragraph{UE} Mobile equipment plus UICC (with the USIM application). LTE UEs are classified
by \emph{categories} specifying peak rates, MIMO layers and buffer sizes.

\paragraph{eNB} The only node of the E-UTRAN \cite{ts36300}. It runs the full radio stack
(PHY, MAC, RLC, PDCP, RRC), schedules uplink and downlink every millisecond, performs admission
control, handover decisions, header compression and ciphering of user data, routes user data to
the S-GW and selects the MME for a UE at attach. eNBs connect to each other over X2 and to the
core over S1.

\paragraph{MME} The control-plane brain of the EPC: terminates NAS signalling and its security,
authenticates the UE with vectors from the HSS, manages tracking areas and paging, selects
S-GW and P-GW, and sets up, modifies and releases bearers through the S-GW. It is the subject of
Lecture~12.

\paragraph{S-GW} Terminates S1-U and is the \emph{local mobility anchor}: during inter-eNB
handover only the S1-U tunnel changes. It buffers downlink data for idle UEs and triggers paging
through the MME, and it anchors mobility to 2G/3G.

\paragraph{P-GW} The \emph{IP anchor}: allocates the UE's IP address, is the point of attachment
to a packet data network (identified by an APN), enforces policy and charging (PCEF), performs
deep packet inspection and lawful interception. A UE may have connections to several P-GWs for
different APNs (e.g.\ \texttt{internet} and \texttt{ims}). Lecture~13 studies both gateways.

\paragraph{HSS} The subscriber database, evolved from the HLR: subscription profile, AuC
functions, current MME. Lecture~11.

\paragraph{PCRF} The policy and charging rules function \cite{ts23203}: decides QoS and charging
rules per service data flow and pushes them to the P-GW over Gx; application functions such as the
IMS P-CSCF request resources over Rx.

\subsection{Reference points}
\begin{table}[htbp]
\centering\small
\caption{Main EPS reference points \cite{ts23401}.}\label{tab:ref}
\rowcolors{2}{acnmist}{white}
\begin{tabularx}{\linewidth}{@{}l l l Y@{}}
\toprule
\rowcolor{white}Reference point & Between & Protocol & Purpose \\
\midrule
LTE-Uu  & UE -- eNB     & RRC, PDCP/RLC/MAC/PHY & Radio interface \\
S1-MME  & eNB -- MME    & S1AP over SCTP        & UE context, bearers, paging, handover \\
S1-U    & eNB -- S-GW   & GTP-U over UDP/IP     & User data per bearer \\
X2      & eNB -- eNB    & X2AP / GTP-U          & Handover, data forwarding, interference coordination \\
S11     & MME -- S-GW   & GTPv2-C               & Session and bearer management \\
S5/S8   & S-GW -- P-GW  & GTPv2-C + GTP-U       & S5 within a network, S8 between home and visited \\
S6a     & MME -- HSS    & Diameter              & Authentication vectors, location update, subscription \\
SGi     & P-GW -- PDN   & IP                    & To the Internet, IMS or corporate networks \\
Gx      & PCRF -- P-GW  & Diameter              & Policy and charging rules \\
Rx      & AF -- PCRF    & Diameter              & Application requests for QoS (e.g.\ VoLTE) \\
Gy/Gz   & P-GW -- OCS/OFCS & Diameter / file    & Online / offline charging \\
\bottomrule
\end{tabularx}
\end{table}

Two patterns stand out. First, every control interface of the core is either GTPv2-C (between
nodes that manage tunnels) or Diameter (towards databases and policy) -- SS7 and MAP have
disappeared from the EPC. Second, S8 is the \emph{roaming} variant of S5: with \emph{home
routing} the visited S-GW tunnels traffic back to the home P-GW, so the home operator keeps
control of IP addressing and charging.

\subsection{Voice in an all-IP system}
Without a CS domain, LTE offers three voice solutions:
\begin{itemize}
  \item \textbf{VoLTE} -- voice as an IMS service \cite{ts23228}. The UE registers with IMS over a
        default bearer to the \texttt{ims} APN (QCI~5 for SIP signalling). At call set-up the
        P-CSCF asks the PCRF over Rx for resources; the PCRF instructs the P-GW over Gx to create a
        dedicated GBR bearer with QCI~1 for the speech packets (Lecture~10).
  \item \textbf{CS fallback (CSFB)} -- the UE is registered in the 2G/3G CS domain via the MME
        (SGs interface); for a call it is redirected to 2G/3G \cite{ts23272}. Simple, but call
        set-up is slower and data speed drops during the call.
  \item \textbf{SRVCC} -- single radio voice call continuity: an ongoing VoLTE call is handed
        over to the CS domain when the UE leaves LTE coverage \cite{ts23216}.
\end{itemize}

\subsection{Control and user plane separation (CUPS)}
In Release~8 the S-GW and P-GW each contained both control logic (GTP-C, policy, IP allocation)
and packet forwarding. Release~14 split them \cite{ts23214}: SGW-C and PGW-C keep the control
logic and program SGW-U and PGW-U through the \emph{Sx} interface with the Packet Forwarding
Control Protocol (PFCP) \cite{ts29244}. The user planes can then be scaled independently and
placed close to the radio for low latency (edge computing), while control remains centralised.
CUPS is the direct ancestor of the 5G SMF/UPF split (Lectures~13 and~19).

% ---------------------------------------------------------------------
\section{Mathematical models}

\subsection{Downlink peak rate from the resource grid}\label{sec:peak}
\paragraph{Problem} Where does the ``300\,Mbit/s'' of a Category-5 LTE device come from?
\paragraph{System model} Bandwidth 20\,MHz, normal cyclic prefix, FDD. The resource grid
(Lecture~7) has $N_{\mathrm{RB}} = 100$ resource blocks of 12 subcarriers at $\Delta f = 15$\,kHz,
and 14 OFDM symbols per 1\,ms subframe, i.e.\ $12\times14 = 168$ resource elements (REs) per RB
per subframe. Assumptions: one OFDM symbol for the control region, four cell-specific
reference-signal ports, 64QAM ($Q_m = 6$), $\nu = 4$ MIMO layers, and we neglect the
synchronisation and broadcast channels, which occupy a few subframes only.
\paragraph{Derivation} Four CRS ports use 24 REs per RB pair, four of which lie in the control
symbol; the data REs per RB per subframe are therefore
\begin{equation}
  N_{\mathrm{RE}}^{\mathrm{data}} = 168 - \underbrace{12}_{\text{control}} - \underbrace{20}_{\text{CRS}} = 136 .
\end{equation}
The peak rate is
\begin{equation}\label{eq:peak}
  R_{\mathrm{peak}} = \nu\, Q_m\, r_\mathrm{c}\,\frac{N_{\mathrm{RB}}\,N_{\mathrm{RE}}^{\mathrm{data}}}{T_{\mathrm{sf}}}
  = 4\times6\times r_\mathrm{c}\times\frac{100\times136}{\qty{1}{\milli\second}}
  = 326.4\,r_\mathrm{c}\ \mathrm{Mbit/s},
\end{equation}
where $r_\mathrm{c}$ is the channel-code rate and $T_{\mathrm{sf}}$ the subframe duration. With the
highest standardised transport block ($r_\mathrm{c}\approx0.92$) this gives
$\approx\qty{300}{\mega\bit\per\second}$, consistent with the 299.6\,Mbit/s of UE Category~5.
\paragraph{Interpretation} The reference signals and control region cost about 19\% of the REs;
the cyclic prefix (6.7\% of time, \cref{exr:cp}) is already built into the 14-symbol subframe.
Peak rates need four receive antennas, excellent SINR and an empty cell -- they are an upper bound
of the air interface, not a user experience.

\subsection{User-plane latency budget}\label{sec:latency}
\paragraph{Problem} Can LTE meet the 5\,ms one-way user-plane latency target?
\paragraph{Model} One-way latency from the arrival of a packet at the eNB's PDCP to its delivery
at the UE's IP layer, unloaded cell, decomposed as
\begin{equation}\label{eq:latency}
  D_{\mathrm{UP}} = \underbrace{\tfrac{1}{2}T_{\mathrm{TTI}}}_{\text{alignment}}
  + T_{\mathrm{TTI}} + D_{\mathrm{eNB}} + D_{\mathrm{UE}} + p_{\mathrm{re}}\,T_{\mathrm{HARQ}}
  + D_{\mathrm{tr}},
\end{equation}
where the alignment term is the mean wait for the next TTI boundary, $D_{\mathrm{eNB}}$ and
$D_{\mathrm{UE}}$ are processing delays, $p_{\mathrm{re}}$ is the probability of one HARQ
retransmission, $T_{\mathrm{HARQ}} = \qty{8}{\milli\second}$ the FDD HARQ round trip and
$D_{\mathrm{tr}}$ the residual transport delay.
\paragraph{Numerical example} With $T_{\mathrm{TTI}} = 1$\,ms, $D_{\mathrm{eNB}} = D_{\mathrm{UE}} = 1$\,ms,
$p_{\mathrm{re}} = 0.1$ and $D_{\mathrm{tr}} = 0.5$\,ms (illustrative values, \cref{fig:latency}):
$D_{\mathrm{UP}} = 0.5 + 1 + 1 + 1 + 0.8 + 0.5 = \qty{4.8}{\milli\second}$.
\paragraph{Interpretation and limits} The 1\,ms TTI and the 8\,ms HARQ loop dominate: these are the
levers that 5G pulls (shorter slots, mini-slots, faster HARQ, Lecture~14). The model ignores
queueing (zero load), scheduling-request delay in the uplink (a UE without a grant must first ask
for one, adding several milliseconds) and core and Internet delay. Measured end-to-end ping times
in LTE networks are therefore typically tens of milliseconds even though the radio target is 5\,ms.

\begin{acnfigure}
  \acnfig{lecture-05/latency-budget}
  \caption{Illustrative one-way downlink user-plane latency budget of \cref{eq:latency} (4.8\,ms).}
  \label{fig:latency}
\end{acnfigure}

\subsection{Encapsulation overhead}\label{sec:overhead}
Every user packet crossing S1-U or S5 is carried in GTP-U/UDP/IPv4, adding
$h = 8 + 8 + 20 = 36$ bytes. For a payload of $L$ bytes the transport efficiency is
\begin{equation}\label{eq:eff}
  \eta_{\mathrm{tr}} = \frac{L}{L + h}.
\end{equation}
A 1400-byte web packet has $\eta_{\mathrm{tr}} = 97.5\%$; a 60-byte VoLTE packet (speech
frame plus RTP/UDP/IP headers; header compression acts only on the radio link) only 62.5\%. Backhaul
dimensioning for voice must account for this (\cref{exr:volte}); IPv6 transport adds a further
20 bytes.

% ---------------------------------------------------------------------
\section{Network architecture and diagrams}
\Cref{fig:arch} shows the EPS. Compared with UMTS (Lecture~3, Figure~3) the RAN has one node type,
the core has no CS domain, and the user plane crosses only two core nodes (S-GW and P-GW),
both of which can be co-located. Control (dashed) and user (solid teal) planes use different
interfaces everywhere except the radio. \Cref{fig:stack} shows the user-plane stack: radio
protocols end in the eNB, and from there GTP-U tunnels carry the user's IP packets.

\begin{acnfigure}
  \resizebox{\linewidth}{!}{\acnfig{lecture-05/lte-architecture}}
  \caption{The Evolved Packet System. Dashed blue: control plane; solid teal: user plane.}
  \label{fig:arch}
\end{acnfigure}

\begin{acnfigure}
  \resizebox{\linewidth}{!}{\acnfig{lecture-05/up-stack}}
  \caption{LTE user-plane protocol stack. The UE's IP packet is carried by PDCP/RLC/MAC over the
  radio and by GTP-U tunnels on S1-U and S5/S8; the P-GW is the IP anchor.}
  \label{fig:stack}
\end{acnfigure}

\begin{acnfigure}
  \acnfig{lecture-05/cups}
  \caption{Control and user plane separation of the EPC gateways (Release~14).}
  \label{fig:cups}
\end{acnfigure}

% ---------------------------------------------------------------------
\section{Step-by-step procedures}
\begin{procedurebox}{LTE initial attach (overview)}\label{proc:attach}
\begin{steps}
  \item The UE performs random access and sets up an RRC connection with the eNB (Lectures~6 and~8).
  \item It sends a NAS \msg{Attach Request} with a \msg{PDN Connectivity Request}; the eNB selects an
        MME and forwards it in S1AP \msg{Initial UE Message}.
  \item The MME authenticates the UE with EPS-AKA (vectors from the HSS over S6a), activates NAS
        security, and registers itself in the HSS (\msg{Update Location}), receiving the subscription.
  \item The MME selects an S-GW and a P-GW and sends GTPv2-C \msg{Create Session Request}; the P-GW
        allocates an IP address, obtains policy from the PCRF and creates the default bearer.
  \item The MME sends S1AP \msg{Initial Context Setup Request} with \msg{Attach Accept}; the eNB
        activates AS security and configures the radio bearer with \msg{RRC Connection Reconfiguration}.
  \item The UE answers \msg{Attach Complete}; the MME sends \msg{Modify Bearer Request} so that the
        S-GW learns the eNB's tunnel endpoint. Uplink and downlink data can flow.
\end{steps}
\end{procedurebox}
The detailed message flow, with every information element that matters, is built up in
Lectures~9, 11 and~12. Note already that one procedure delivers registration, security
\emph{and} an IP address -- three separate GPRS procedures merged into one.

% ---------------------------------------------------------------------
\section{Worked examples}

\begin{examplebox}{Peak rate of a 10\,MHz cell}\label{ex:peak10}
\textbf{Problem.} A 10\,MHz LTE carrier (50 RBs) uses two CRS ports, two control symbols and
$2\times2$ MIMO with 64QAM. Estimate the peak rate for $r_\mathrm{c} = 0.9$.

\textbf{Solution.} Two CRS ports use 16 REs per RB pair, of which 4 lie in the first (control)
symbol; with two control symbols, $N_{\mathrm{RE}}^{\mathrm{data}} = 168 - 24 - 12 = 132$.
$R = 2\times6\times0.9\times50\times132/10^{-3} = \qty{71.3}{\mega\bit\per\second}$.

\textbf{Interpretation.} Half the bandwidth and half the layers give roughly a quarter of the
Category-5 figure; the control region and RS overhead are similar ($\approx21\%$).
\end{examplebox}

\begin{examplebox}{Hops in the user plane}\label{ex:hops}
\textbf{Problem.} Count the network nodes that process a user packet between the radio and the
Internet in UMTS R99 and in LTE, and estimate the saving if each node adds 0.5\,ms of processing
and queuing.

\textbf{Solution.} UMTS: NodeB, RNC, SGSN, GGSN (4 nodes). LTE: eNB, S-GW, P-GW (3 nodes; 2 if
S-GW and P-GW are combined). Saving: 0.5--1\,ms one way, 1--2\,ms round trip.

\textbf{Interpretation.} The node count matters less than the RNC's long TTIs and slow
scheduling, which the eNB replaces by 1\,ms scheduling -- the main reason LTE's latency is a
fraction of R99's.
\end{examplebox}

\begin{examplebox}{VoLTE capacity of a cell}\label{ex:volte}
\textbf{Problem.} A VoLTE call needs on average 2 RBs every 20\,ms in each direction (one
speech frame per 20\,ms, robust MCS). How many calls can a 10\,MHz cell carry if 70\% of RBs may
be used for voice?

\textbf{Solution.} Per call: 2 RBs per 20 subframes $=$ 0.1 RB per subframe. Available: $0.7\times50
= 35$ RBs per subframe. Calls $= 35/0.1 = 350$ per direction (in practice the PDCCH control channel
limits this earlier, which is why \emph{semi-persistent scheduling} is used for VoLTE).

\textbf{Interpretation.} Compare with about 45 WCDMA voice users per 5\,MHz (Lecture~3): packet
voice with header compression and scheduling is far more spectrum-efficient than circuit voice.
\end{examplebox}

% ---------------------------------------------------------------------
\section{Technology comparison}
\begin{table}[htbp]
\centering\small
\caption{3G (UMTS/HSPA) versus 4G (LTE/EPC).}\label{tab:3g4g}
\rowcolors{2}{acnmist}{white}
\begin{tabularx}{\linewidth}{@{}l Y Y@{}}
\toprule
\rowcolor{white} & UMTS / HSPA & LTE / EPC \\
\midrule
Radio access          & WCDMA, 5\,MHz                     & OFDMA DL / SC-FDMA UL, 1.4--20\,MHz \\
RAN nodes             & NodeB + RNC                        & eNB only (flat) \\
Handover              & Soft + hard                        & Hard (with X2 forwarding) \\
Core domains          & CS + PS                            & PS only \\
Voice                 & CS                                 & VoLTE (IMS); CSFB/SRVCC for interworking \\
Core signalling       & MAP (SS7), GTPv1-C, RANAP          & Diameter, GTPv2-C, S1AP \\
TTI                   & 2\,ms (HSPA), 10--80\,ms (R99)     & 1\,ms \\
Always-on IP          & No (PDP context on demand)          & Yes (default bearer at attach) \\
Peak DL (typical device) & 14--42\,Mbit/s                  & 150--300\,Mbit/s (Rel-8) \\
\bottomrule
\end{tabularx}
\end{table}

\paragraph{Trade-offs} The flat, all-IP design simplified the network but moved complexity
elsewhere: QoS for voice must now be engineered with bearers and policy (PCRF), the eNB carries
more functions and state, mobility without soft handover relies on fast hard handovers, and the
transition period needed CSFB and SRVCC. The control/user plane separation already present in
the EPC (MME versus gateways) and completed by CUPS foreshadows the service-based 5G core.

% ---------------------------------------------------------------------
\section{Key takeaways}
\begin{takeawaybox}
\begin{itemize}
  \item LTE requirements: 100/50\,Mbit/s peak in 20\,MHz, <5\,ms user-plane latency, scalable
        bandwidth, packet-only.
  \item EPS = E-UTRAN (eNB only) + EPC (MME, S-GW, P-GW, HSS, PCRF).
  \item Control interfaces use S1AP, GTPv2-C and Diameter; user data uses GTP-U on S1-U and S5/S8.
  \item S-GW = mobility anchor; P-GW = IP anchor and policy enforcement; MME = signalling only.
  \item A default EPS bearer is created at attach: always-on IP.
  \item Peak rate $\approx\nu Q_m r_\mathrm{c} N_{\mathrm{RB}} N_{\mathrm{RE}}^{\mathrm{data}}/T_{\mathrm{sf}}$;
        the 1\,ms TTI and 8\,ms HARQ loop dominate the radio latency.
  \item Voice: VoLTE over IMS with a QCI~1 bearer; CSFB and SRVCC for interworking with 2G/3G.
  \item CUPS splits the gateways into C and U parts controlled with PFCP -- the bridge to 5G.
\end{itemize}
\end{takeawaybox}

% ---------------------------------------------------------------------
\section{Exercises}

\begin{exercisebox}{1}{Why two gateways?}\label{exr:twogw}
Explain why the EPC has both an S-GW and a P-GW. Give one deployment where they are combined and
one where they must be separate.
\end{exercisebox}

\begin{exercisebox}{1}{Where did the RNC go?}\label{exr:rnc}
For each of the following RNC functions state where it lives in LTE: RRC termination, ciphering of
user data, handover decision, outer-loop power control, mobility anchoring towards the core.
\end{exercisebox}

\begin{exercisebox}{2}{Cyclic-prefix overhead}\label{exr:cp}
In LTE with normal CP, a 0.5\,ms slot contains 7 OFDM symbols of $2048\,T_\mathrm{s}$ with
$T_\mathrm{s} = 1/(15000\times2048)$\,s; the first CP lasts $160\,T_\mathrm{s}$ and the others
$144\,T_\mathrm{s}$. Compute the CP durations in microseconds and the fraction of time spent on
CP. What is the maximum delay spread tolerated without inter-symbol interference?
\end{exercisebox}

\begin{exercisebox}{2}{VoLTE backhaul}\label{exr:volte}
A VoLTE packet carries a 33-byte speech frame every 20\,ms with 12 bytes RTP, 8 bytes UDP and
20 bytes IPv4. On S1-U it is encapsulated in GTP-U/UDP/IPv4. Compute the bit rate per call on the
backhaul and the transport efficiency \cref{eq:eff}. How many calls fit in a 100\,Mbit/s link
loaded at 70\%?
\end{exercisebox}

\begin{exercisebox}{3}{Uplink latency}\label{exr:ullat}
Extend \cref{eq:latency} to the uplink of a UE that has no grant: it must wait for its next
scheduling-request (SR) opportunity (period $T_{\mathrm{SR}}$), send the SR, receive a grant after
$k_1$\,ms and transmit $k_2$\,ms after the grant. Derive the mean latency and evaluate it for
$T_{\mathrm{SR}} = 10$\,ms, $k_1 = 3$\,ms, $k_2 = 4$\,ms, other values as in \cref{sec:latency}.
\end{exercisebox}

\begin{exercisebox}{3}{Home routing versus local breakout}\label{exr:roaming}
A subscriber from Tehran roams in Istanbul and browses a server in Istanbul. Compare the user-plane
path and the round-trip propagation delay for home-routed roaming (S8 to the home P-GW) and local
breakout (visited P-GW). Assume 2000\,km between the cities, fibre at $2\times10^8$\,m/s and a route
factor of 1.5. What does each option mean for charging and lawful interception?
\end{exercisebox}

\subsection*{Solutions}
\addcontentsline{toc}{subsection}{Solutions}

\begin{solutionbox}{\ref{exr:twogw}}
They anchor different things. The S-GW anchors mobility between eNBs (and to 2G/3G) and buffers
for idle UEs; there is one per UE at a time and it may change with location. The P-GW anchors the
IP address and the connection to a PDN and enforces policy; it must not change during a session.
Combined: a small national operator, one site, no roaming. Separate: roaming with home routing
(visited S-GW, home P-GW over S8), or a UE with an Internet PDN on a local P-GW and an IMS PDN
on a central P-GW.
\end{solutionbox}

\begin{solutionbox}{\ref{exr:rnc}}
RRC termination: eNB. User-data ciphering: eNB (PDCP). Handover decision: source eNB. Outer-loop
power control: not needed in the same form -- LTE uses scheduled, orthogonal resources; uplink power
control (open loop plus closed-loop corrections) is in the eNB. Mobility anchoring towards the
core: S-GW (the RNC's Iu anchor role).
\end{solutionbox}

\begin{solutionbox}{\ref{exr:cp}}
$T_\mathrm{s} = 32.55$\,ns. First CP: $160\,T_\mathrm{s} = \qty{5.21}{\micro\second}$; others:
$144\,T_\mathrm{s} = \qty{4.69}{\micro\second}$; useful symbol $2048\,T_\mathrm{s} = \qty{66.67}{\micro\second}$.
CP fraction: $(160+6\times144)/15360 = 6.67\%$. Delay spreads up to about $\qty{4.7}{\micro\second}$
(about 1.4\,km of path difference) are absorbed; the extended CP ($512\,T_\mathrm{s} =
\qty{16.7}{\micro\second}$) serves large cells and broadcast.
\end{solutionbox}

\begin{solutionbox}{\ref{exr:volte}}
IP packet: $33+12+8+20 = 73$ bytes; on S1-U: $73+36 = 109$ bytes every 20\,ms $=
\qty{43.6}{\kilo\bit\per\second}$ (plus Ethernet framing in practice). Efficiency relative to the
speech payload: $33/109 = 30\%$; relative to the user IP packet \cref{eq:eff}: $73/109 = 67\%$.
Calls: $0.7\times100\times10^6/43\,600 \approx 1600$ per direction.
\end{solutionbox}

\begin{solutionbox}{\ref{exr:ullat}}
Mean wait for an SR opportunity $T_{\mathrm{SR}}/2$, then SR transmission (1\,ms), grant after $k_1$,
data $k_2$ after the grant, then the downlink-like terms:
$D_{\mathrm{UL}} = T_{\mathrm{SR}}/2 + T_{\mathrm{TTI}} + k_1 + k_2 + T_{\mathrm{TTI}} + D_{\mathrm{eNB}} +
p_{\mathrm{re}}T_{\mathrm{HARQ}} + D_{\mathrm{tr}} = 5 + 1 + 3 + 4 + 1 + 1 + 0.8 + 0.5 = \qty{16.3}{\milli\second}$.
The scheduling-request cycle dominates; pre-scheduling or short SR periods reduce it, and 5G adds
configured grants that remove it.
\end{solutionbox}

\begin{solutionbox}{\ref{exr:roaming}}
Home routing: Istanbul UE $\to$ visited S-GW $\to$ S8 to Tehran P-GW $\to$ back to the Istanbul
server: two traversals of $1.5\times2000 = 3000$\,km each way, i.e.\ 12\,000\,km per round trip, or
60\,ms of propagation. Local breakout: essentially local, $\ll1$\,ms. Home routing keeps charging,
policy and lawful interception in the home network (simple for the home operator); local breakout
requires the visited network to apply the home operator's policy (via S9 / roaming agreements) and
to perform interception locally.
\end{solutionbox}

% ---------------------------------------------------------------------
\section{Further reading}
The E-UTRAN overall description is \cite{ts36300} and the EPC procedures \cite{ts23401}; the LTE
requirements are in \cite{tr25913}. Policy and charging control is specified in \cite{ts23203},
CUPS in \cite{ts23214} with PFCP in \cite{ts29244}, CSFB in \cite{ts23272} and SRVCC in
\cite{ts23216}. Sesia et al.\ \cite{sesia2011lte} and Olsson et al.\ \cite{olsson2013epc} are the
recommended textbooks for the radio and core sides respectively.

\printbibliography[heading=subbibliography,title={\ACNstrRefs}]
\end{document}
